\documentclass[10pt,a4paper]{amsart}
\usepackage[margin=19mm]{geometry}
\usepackage{amsmath,amssymb,mathtools}
\usepackage[scr=rsfs,cal=euler]{mathalfa}
\usepackage{xfrac}
\usepackage[unicode,hidelinks,bookmarks=false]{hyperref}
\newcommand{\ie}{i.e.\ }
\newcommand{\cf}{cf.\ }
\newcommand{\resp}{respectively\ }
\newcommand{\Subsection}{\subsection}
\providecommand{\nobreakdash}{\nobreak\hbox{-}}
\setlength{\emergencystretch}{2em}
\multlinegap0pt
\usepackage{amssymb}
\usepackage{dsfont}
\newtheorem{lemma}{Lemma}
\title{Commuting maps of inflated algebras}
\author{Hongyu Jia, Zhankui Xiao}
\date{Source: arXiv:2605.12849v1 (CC BY 4.0)}
\begin{document}
\maketitle

\noindent\emph{Source and licence.} Commuting maps of inflated algebras. Version arXiv:2605.12849v1 (CC BY 4.0), distributed by arXiv under the Creative Commons Attribution 4.0 International licence, http://creativecommons.org/licenses/by/4.0/, as recorded in the arXiv licence metadata for this exact version and shown on the abstract page https://arxiv.org/abs/2605.12849v1. Full licence notice: \texttt{licenses/arxiv-2605.12849-CC-BY-4.0.txt}.\par\smallskip

\noindent\textit{Editorial scope.} Counted unit: the article's lemma l22 (source0.tex lines 287--294). The article's complete original proof of it is delivered with it (source0.tex lines 295--491, one \texttt{proof} environment of 197 lines). No mathematics is written, repaired or re-derived: the statement and the proof are the article's own lines, in the article's own notation, and no retained line is substituted or re-flowed.  Retained uncounted as the source-local context the proof needs: lines 277--279.  Nothing was omitted from any retained span, so the package carries no omission record.  No retained line contains a citation, so the wrapper carries no bibliography and no bibliography entry.  No retained line contains a figure environment, an image-inclusion command or drawing code, so no image is delivered and the package declares no asset dependency.  Editorial formatting only, in addition: an AMS \texttt{amsart} preamble that restates the article's own declarations and macros used by the retained lines, namely the environments \texttt{lemma} on separate counters, together with the standard packages those lines need.  The retained environments print in this document as Lemma 1 in order of appearance, whereas the article prints the same items with the numbering of its own full document.  The complete native source of the article as distributed in the arXiv e-print for this exact version is bundled unchanged as \texttt{sources/200457/source0.tex}.
\begin{lemma}\label{l21}
Let $\mathcal{A}$ be a $K$-algebra with a $K$-basis $\Gamma$. Then a $K$-linear map $\theta: \mathcal{A} \to \mathcal{A}$ is a commuting map if and only if $[\theta(x), y] = [x, \theta(y)]$ for all $x, y\in \Gamma$.
\end{lemma}

\begin{lemma}\label{l22}
The commuting map $\theta$ satisfies:
 \begin{align}
 \theta(e_{ii}) &= \sum_{x \in \mathcal{J}} C_{xx}^{ii} e_{xx} + \sum_{x,y \in \hat{\mathcal{J}}} C_{xy}^{ii} e_{xy}, \quad i   \in \mathbf{n},  \label{eq1}\\
 \theta(e_{ij}) &= \sum_{x \in \mathcal{J}} C_{xx}^{ij} e_{xx} + C_{ij}^{ij} e_{ij} + 
 \sum_{\substack{(x,y) \in \hat{\mathcal{J}}\times \hat{\mathcal{J}}\setminus \{(i,j)\}}} C_{xy}^{ij} e_{xy}, \quad i \neq j \in \mathbf{n}.\label{eq2}
 \end{align}
\end{lemma}

\begin{proof}
Without loss of generality, we assume that $|\mathbf{n}| \geq 2$ and $|\mathcal{J}| \geq 1$. In order to determine the coefficients $C_{xy}^{ii}$, we consider in two cases.

\textbf{Case 1.1.} $i \in \mathcal{J}$.
 Since $\theta$ is a commuting map,  we get $[\theta(e_{ii}), e_{ii}] = 0$. Let $k \in \mathbf{n} \setminus \{i\}$. Equating the coefficients of $e_{ik}$ and $e_{ki}$, we have
  \begin{align}\label{eq3}
  C_{ik}^{ii} = C_{ki}^{ii} = 0, \quad \text{for all } k \in \mathbf{n} \setminus \{i\}.
  \end{align}
Since $[\theta(e_{ii}), e_{xx}] = [e_{ii}, \theta(e_{xx})]$ for all $x \in \mathcal{J}$ by Lemma \ref{l21}, equating the coefficients of $e_{xy}$ and $e_{yx}$, with $y \neq x$ and $y \neq i$, then we have
 \begin{align}\label{eq4}
 C_{xy}^{ii} = C_{yx}^{ii} = 0, \quad \text{for all } x \in \mathcal{J}, y \in \mathbf{n} \setminus \{i,x\}.
 \end{align}

Combining the equations (\ref{eq3}) and (\ref{eq4}), we obtain
 \begin{align*}
 \theta(e_{ii})=& \sum_{x,y \in \mathcal{J}} C_{xy}^{ii} e_{xy}
 +\sum_{(x,y) \in \mathcal{J} \times \hat{\mathcal{J}}} (C_{xy}^{ii} e_{xy}+C_{yx}^{ii} e_{yx})+\sum_{x,y \in \hat{\mathcal{J}}} C_{xy}^{ii} e_{xy}\\
=&\sum_{x \in \mathcal{J}} C_{xx}^{ii} e_{xx} + \sum_{x,y \in \mathcal{J}, x\neq y} C_{xy}^{ii} e_{xy}
 +\sum_{y\in \hat{\mathcal{J}}} (C_{iy}^{ii} e_{iy}+C_{yi}^{ii} e_{yi})\\
  &+\sum_{x \in \mathcal{J}\setminus \{i\}, y\in \hat{\mathcal{J}}} (C_{xy}^{ii} e_{xy}+C_{yx}^{ii} e_{yx})+\sum_{x,y \in \hat{\mathcal{J}}} C_{xy}^{ii} e_{xy}\\
=& \sum_{x \in \mathcal{J}} C_{xx}^{ii} e_{xx} + \sum_{x, y \in \hat{\mathcal{J}}} C_{xy}^{ii} e_{xy}.
 \end{align*}
This proves (\ref{eq1}) for the case $ i \in \mathcal{J} $.

\textbf{Case 1.2.} $ i \in \hat{\mathcal{J}} $. Since $ [\theta(e_{ii}), e_{xx}] = [e_{ii}, \theta(e_{xx})] = e_{ii} \cdot  \theta(e_{xx}) - \theta(e_{xx}) \cdot  e_{ii} = 0 $ for all $ x \in \mathcal{J} $, we have 
\[
0 = \theta(e_{ii}) \cdot  e_{xx} - e_{xx} \cdot  \theta(e_{ii}) = \theta(e_{ii}) e_{xx} - e_{xx} \theta(e_{ii}).
\]
Then, equating the coefficients of $ e_{xy} $ and $ e_{yx} $, where $ y \in \mathbf{n} \setminus \{x\} $, we obtain 
 \begin{align}\label{eq5}
 C_{xy}^{ii} = C_{yx}^{ii} = 0, \quad \text{for all } x \in \mathcal{J},\, y \in \mathbf{n} \setminus \{x\}.
 \end{align}
It follows from the equation (\ref{eq5}) that
\[
\begin{aligned}
&\theta(e_{ii}) = \sum_{\substack{x,y \in \mathbf{n}}} C_{xy}^{ii} e_{xy} \\
=&\sum_{x,y \in \mathcal{J}} C_{xy}^{ii} e_{xy}
 +\sum_{(x,y) \in \mathcal{J} \times \hat{\mathcal{J}}} (C_{xy}^{ii} e_{xy}+C_{yx}^{ii} e_{yx})+\sum_{x,y \in \hat{\mathcal{J}}} C_{xy}^{ii} e_{xy}\\
=& \sum_{x \in \mathcal{J}} C_{xx}^{ii} e_{xx} + \sum_{x,y \in \mathcal{J}, x\neq y} C_{xy}^{ii} e_{xy}+
\sum_{x \in \mathcal{J}, y\in \hat{\mathcal{J}}} (C_{xy}^{ii} e_{xy}+C_{yx}^{ii} e_{yx})+\sum_{x,y \in \hat{\mathcal{J}}} C_{xy}^{ii} e_{xy}\\
=& \sum_{x \in \mathcal{J}} C_{xx}^{ii} e_{xx} + \sum_{\substack{x , y \in \hat{\mathcal{J}}}} C_{xy}^{ii} e_{xy}.
\end{aligned}
\]
This proves (\ref{eq1}) for the case $ i \in \hat{\mathcal{J}} $.


For all $ i,j \in \mathbf{n} $ with $ i \neq j $, we will prove that the coefficients of $\theta(e_{ij})$ satisfy equation (\ref{eq2}) by dividing into four cases.

\textbf{Case 2.1.} $ i,j \in \mathcal{J} $. Let $ x \in \mathcal{J} $. Since $ \theta $ is a commuting map, we have $ [\theta(e_{ij}), e_{xx}] = [e_{ij}, \theta(e_{xx})] $.
On the one hand, by (\ref{eq1}), 
\[
\begin{aligned} 
\relax[e_{ij},\theta(e_{xx})]&= e_{ij}\cdot \theta(e_{xx})-\theta(e_{xx})\cdot  e_{ij} = e_{ij} ( \sum_{y \in \mathcal{J}} C_{yy}^{xx} e_{yy})-(\sum_{y \in \mathcal{J}} C_{yy}^{xx} e_{yy}) e_{ij} \\
&= ( C_{jj}^{xx} - C_{ii}^{xx} ) e_{ij}.
\end{aligned}
\]
On the other hand, 
$
[\theta(e_{ij}), e_{xx}] = \theta(e_{ij}) \cdot  e_{xx} - e_{xx} \cdot  \theta(e_{ij}) = \theta(e_{ij}) e_{xx} - e_{xx} \theta(e_{ij}).
$
Combining the just above two identities, we get  
 \begin{align}\label{eq6}
 \theta(e_{ij}) e_{xx} - e_{xx} \theta(e_{ij}) = \left( C_{jj}^{xx} - C_{ii}^{xx} \right) e_{ij},   \text{ for all }  i, j,x \in \mathcal{J}.
 \end{align}
If $ x = i $ and $ y \in \mathbf{n} \setminus \{j\} $, then, by equating the coefficients of $ e_{iy} $ and $ e_{ij} $, we have 
 \begin{align}\label{eq7}
 C_{iy}^{ij} = 0, \quad \text{for all } i \neq j \in \mathcal{J} \text{ and } y \in \mathbf{n} \setminus \{j\}.
 \end{align}
and 
 \begin{align}\label{eq8}
C_{ij}^{ij} = C_{ii}^{ii} - C_{jj}^{ii}, \quad \text{for all } i \neq j \in \mathcal{J}. 
\end{align}
If $ x = j $ in (\ref{eq6}), for all $ y \in \mathbf{n} \setminus \{i\} $, then  by equating the coefficients of $ e_{yj} $ and $ e_{ij} $, we have 
 \begin{align}\label{eq9}
C_{yj}^{ij} = 0, \quad \text{for all } i \neq j \in \mathcal{J} \text{ and } y \in \mathbf{n} \setminus \{i\}. 
\end{align}
and 
 \begin{align}\label{eq10}
C_{ij}^{ij} = C_{jj}^{jj} - C_{ii}^{jj}, \quad \text{for all }  i \neq j \in \mathcal{J}. 
 \end{align}
If $ x \neq i $ in (\ref{eq6}), then $ e_{xx} \theta(e_{ij}) (I_n - e_{xx}) = 0 $, we have
 \begin{align}\label{eq11}
C_{xy}^{ij} = 0, \quad x \in \mathcal{J} \setminus \{i\}, \, y \in \mathbf{n} \setminus \{x\}.
\end{align}
If $ x \neq j $ in (\ref{eq6}), then $ (I_n - e_{xx}) \theta(e_{ij}) e_{xx} = 0 $, we have
 \begin{align}\label{eq12}
C_{yx}^{ij} = 0, \quad x \in \mathcal{J} \setminus \{j\}, \, y \in \mathbf{n} \setminus \{x\}.
\end{align}
Combining the equations (\ref{eq7}), (\ref{eq9}), (\ref{eq11}) and (\ref{eq12}), we can express $\theta(e_{ij})$ as follows
\[
\begin{aligned}
&\theta(e_{ij})=\sum_{\substack{x,y \in \mathbf{n}}} C_{xy}^{ij} e_{xy}\\
=& \sum_{x,y \in \mathcal{J}} C_{xy}^{ij} e_{xy} + \sum_{\substack{x \in \mathcal{J} \\ y \in \hat{\mathcal{J}}}} C_{xy}^{ij} e_{xy} + \sum_{\substack{x \in \hat{\mathcal{J}} \\ y \in \mathcal{J}}} C_{xy}^{ij} e_{xy} + \sum_{\substack{x,y \in \hat{\mathcal{J}}}} C_{xy}^{ij} e_{xy} \\
=& \sum_{x \in \mathcal{J}} C_{xx}^{ij} e_{xx} +\sum_{\substack{x \in \mathcal{J}\setminus \{i\} \\y \in \mathcal{J}\setminus \{j\}, x\neq y }} C_{xy}^{ij} e_{xy} + \sum_{\substack{x \in \mathcal{J}\setminus \{i\} \\y \in \hat{\mathcal{J}} }} C_{xy}^{ij} e_{xy} +
\sum_{\substack{x \in \hat{\mathcal{J}}\\y \in \mathcal{J}\setminus \{j\} }} C_{xy}^{ij} e_{xy} \\
&+\sum_{\substack{x \in \mathcal{J}\setminus \{i\}  }} C_{xj}^{ij} e_{xj}
\sum_{\substack{y \in \mathcal{J}\setminus \{j\} }} C_{iy}^{ij} e_{iy}+C_{ij}^{ij} e_{ij}
+ \sum_{\substack{x,y \in \hat{\mathcal{J}}}} C_{xy}^{ij} e_{xy} \\
=& \sum_{x \in \mathcal{J}} C_{xx}^{ij} e_{xx} + C_{ij}^{ij} e_{ij} + 
 \sum_{\substack{(x,y) \in \hat{\mathcal{J}}\times \hat{\mathcal{J}}}} C_{xy}^{ij} e_{xy}.
\end{aligned}
\]
Thus we prove (\ref{eq2}) for the case $ i \neq j \in \mathcal{J} $.

\textbf{Case 2.2.} $ i \in \mathcal{J},~j \in \hat{\mathcal{J}} $. Let $ x \in \mathcal{J} $. Since $ \theta $ is a commuting map, $ [\theta(e_{ij}), e_{xx}] = [e_{ij}, \theta(e_{xx})] $.
On the one hand, by (\ref{eq1}),
$
[e_{ij}, \theta(e_{xx})] = e_{ij} \cdot  \theta(e_{xx}) - \theta(e_{xx}) \cdot  e_{ij} = -C_{ii}^{xx} e_{ij}.
$
On the other hand, $ [\theta(e_{ij}), e_{xx}] = \theta(e_{ij}) e_{xx} - e_{xx} \theta(e_{ij}) $.
It follows that
 \begin{align}\label{eq13}
\theta(e_{ij}) e_{xx} - e_{xx} \theta(e_{ij}) = -C_{ii}^{xx} e_{ij} .
 \end{align}
If $ x = i $ in  (\ref{eq13}), then  by equating the coefficients of $ e_{iy} $ and $ e_{ij} $, we have
 \begin{align}\label{eq14}
C_{iy}^{ij} = 0, \quad \text{for all } y \in \mathbf{n} \setminus \{i, j\},
 \end{align}
and
 \begin{align}\label{eq15}
C_{ij}^{ij} = C_{ii}^{ii},\quad \text{for all } i \in \mathcal{J}, j \in \hat{\mathcal{J}}. 
 \end{align}
If $ x \neq i $ in (\ref{eq13}), then  $ e_{xx} \theta(e_{ij}) (I_n - e_{xx}) = 0 $ and hence
 \begin{align}\label{eq16}
C_{xy}^{ij} = 0, \quad \text{for all } x \in \mathcal{J} \setminus \{i\}, y \in \mathbf{n} \setminus \{x\}.
 \end{align}
Since $ x \in \mathcal{J}$ and $ j \in \hat{\mathcal{J}} $ imply  $ x \neq j $, we have $ (I_n - e_{xx}) \theta(e_{ij}) e_{xx} = 0 $ for all $x\in \mathcal{J}$ by (\ref{eq13}). This leads to
 \begin{align}\label{eq17}
C_{yx}^{ij} = 0, \quad \text{for all } x \in \mathcal{J}, y \in \mathbf{n} \setminus \{x\}.
 \end{align}
Now combining the equations (\ref{eq14}), (\ref{eq16}) and (\ref{eq17}), we obtain
\[
\begin{aligned}
\theta(e_{ij}) 
=& \sum_{\substack{x,y \in \mathbf{n}}} C_{xy}^{ij} e_{xy}=\sum_{\substack{x \in \mathcal{J} \\ y \in \mathbf{n}}} C_{xy}^{ij} e_{xy} + \sum_{\substack{x \in \hat{\mathcal{J}} \\ y \in \mathcal{J}}} C_{xy}^{ij} e_{xy} + \sum_{\substack{x,y \in \hat{\mathcal{J}}}} C_{xy}^{ij} e_{xy} \\
=& \sum_{x \in \mathcal{J}} C_{xx}^{ij} e_{xx} + \sum_{\substack{y \in \mathbf{n}\setminus \{i,j\} }} C_{iy}^{ij} e_{iy}+C_{ij}^{ij} e_{ij}
+\sum_{\substack{x \in \mathcal{J}\setminus \{i\} \\y \in \mathbf{n}\setminus \{x\} }} C_{xy}^{ij} e_{xy} \\
&+\sum_{\substack{x \in \hat{\mathcal{J}}\\y \in \mathcal{J} }} C_{xy}^{ij} e_{xy}+ \sum_{\substack{x,y \in \hat{\mathcal{J}}}} C_{xy}^{ij} e_{xy} \\
=& \sum_{x \in \mathcal{J}} C_{xx}^{ij} e_{xx} + C_{ij}^{ij} e_{ij} + 
 \sum_{\substack{(x,y) \in \hat{\mathcal{J}}\times \hat{\mathcal{J}}}} C_{xy}^{ij} e_{xy}.
\end{aligned}
\]
This proves (\ref{eq2}) for the case $i \in \mathcal{J}, ~j \in \hat{\mathcal{J}}$.

\textbf{Case 2.3.} $i \in \hat{\mathcal{J}},~j \in \mathcal{J}$. Let $x \in \mathcal{J}$. Since $\theta$ is a commuting map, we have $[\theta(e_{ij}), e_{xx}] = [e_{ij}, \theta(e_{xx})]$.
On the one hand, by (\ref{eq1}) and the multiplication of $\mathfrak{M}_n(r)$, $[e_{ij}, \theta(e_{xx})] = e_{ij} \cdot  \theta(e_{xx}) - \theta(e_{xx}) \cdot  e_{ij} = C_{jj}^{xx} e_{ij}$.
On the other hand, $[\theta(e_{ij}), e_{xx}] = \theta(e_{ij}) e_{xx} - e_{xx} \theta(e_{ij})$. It follows that
 \begin{align}\label{eq18}
\theta(e_{ij}) e_{xx} - e_{xx} \theta(e_{ij}) = C_{jj}^{xx} e_{ij}.
 \end{align}
Since $x \in \mathcal{J},~i \in \hat{\mathcal{J}}$ imply $x \neq i$, we obtain $e_{xx} \theta(e_{ij}) (I_n - e_{xx}) = 0$. This leads to
 \begin{align}\label{eq19}
C_{xy}^{ij} = 0, \quad \text{for all } x \in \mathcal{J}, ~y \in \mathbf{n} \setminus \{x\}.  
\end{align}
If $x = j$ in (\ref{eq18}). Then, by equating the coefficients of $e_{yj}$ and $e_{ij}$, $y \in \mathbf{n} \setminus \{i, j\}$, we have
 \begin{align}\label{eq20}
C_{yj}^{ij} = 0, \quad \text{for all } y \in \mathbf{n} \setminus \{i, j\},
 \end{align}
and
 \begin{align}\label{eq21}
C_{ij}^{ij} = C_{jj}^{jj}, \quad \text{for all } i \in \hat{\mathcal{J}}, ~j \in \mathcal{J}. 
 \end{align}
If $x \neq j$ in (\ref{eq18}). Then $(I_n - e_{xx}) \theta(e_{ij}) e_{xx} = 0$, we have
 \begin{align}\label{eq22}
C_{yx}^{ij} = 0, \quad \text{for all } x \in \mathcal{J} \setminus \{j\},~ y \in \mathbf{n} \setminus \{x\} .
\end{align}
Combining the equations (\ref{eq19}), (\ref{eq20}) and (\ref{eq22}), we obtain
\[
\begin{aligned}
\theta(e_{ij})
=& \sum_{\substack{x,y \in \mathbf{n}}} C_{xy}^{ij} e_{xy}
= \sum_{\substack{x \in \mathbf{n} \\ y \in \mathcal{J}}} C_{xy}^{ij} e_{xy}+ \sum_{\substack{x \in \mathcal{J} \\ y \in \hat{\mathcal{J}}}} C_{xy}^{ij} e_{xy} +\sum_{\substack{x,y \in \hat{\mathcal{J}}}} C_{xy}^{ij} e_{xy}\\
=&\sum_{x \in \mathcal{J}} C_{xx}^{ij} e_{xx} + \sum_{\substack{x \in \mathbf{n}\setminus \{i,j\} }} C_{xj}^{ij} e_{xj}+C_{ij}^{ij} e_{ij}+\sum_{\substack{x \in \mathcal{J} \\y \in \mathbf{n}\setminus \{j, x\} }} C_{xy}^{ij} e_{xy} \\
&+\sum_{\substack{x \in \mathcal{J}\\y \in \hat{\mathcal{J}} }} C_{xy}^{ij} e_{xy}+ \sum_{\substack{x,y \in \hat{\mathcal{J}}}} C_{xy}^{ij} e_{xy}  \\
=& \sum_{x \in \mathcal{J}} C_{xx}^{ij} e_{xx} + C_{ij}^{ij} e_{ij} + 
 \sum_{\substack{(x,y) \in \hat{\mathcal{J}}\times \hat{\mathcal{J}}}} C_{xy}^{ij} e_{xy}.
\end{aligned}
\]
This proves (\ref{eq2}) for the case $i \in \hat{\mathcal{J}}, ~j \in \mathcal{J}$.

\textbf{Case 2.4.} $i, j \in \hat{\mathcal{J}}$ and $i \neq j$. Let $x \in \mathcal{J}$. Since $\theta$ is a commuting map, we have $[\theta(e_{ij}), e_{xx}] = [e_{ij}, \theta(e_{xx})] = 0$ and therefore $\theta(e_{ij}) e_{xx} - e_{xx} \theta(e_{ij}) = 0$.
By equating the coefficients of $e_{yx}$ and $e_{xy}$, where $y \in \mathbf{n} \setminus \{x\}$, we have
 \begin{align}\label{eq23}
C_{xy}^{ij} = C_{yx}^{ij} = 0, \quad \text{for all } x \in \mathcal{J},~y \in \mathbf{n} \setminus \{x\}.
\end{align}
It follows from the equation (\ref{eq23}) that
\[
\begin{aligned}
\theta(e_{ij}) =& \sum_{\substack{x,y \in \mathbf{n}}} C_{xy}^{ij} e_{xy} \\
=& \sum_{x \in \mathcal{J}} C_{xx}^{ij} e_{xx} + \sum_{\substack{x \in \hat{\mathcal{J}} \\ y \in \mathcal{J}}} C_{xy}^{ij} e_{xy} + \sum_{\substack{x \in \mathcal{J} \\ y \in \mathbf{n} \setminus \{x\}}} C_{xy}^{ij} e_{xy} + \sum_{\substack{x,y \in \hat{\mathcal{J}}}} C_{xy}^{ij} e_{xy} \\
=& \sum_{x \in \mathcal{J}} C_{xx}^{ij} e_{xx}  + \sum_{\substack{(x,y) \in \hat{\mathcal{J}}\times \hat{\mathcal{J}}}} C_{xy}^{ij} e_{xy}\\
=& \sum_{x \in \mathcal{J}} C_{xx}^{ij} e_{xx} + C_{ij}^{ij} e_{ij} + \sum_{\substack{(x,y) \in \hat{\mathcal{J}}\times \hat{\mathcal{J}}\setminus \{(i,j)\}}} C_{xy}^{ij} e_{xy}.
\end{aligned}
\]
%%%%%%%%%%%%%%%%%%%%%%%%%%%%%%%%%%%%%%%%%%%%%%%%%%%%%%%%%%%%%%%%%%%%%%%%%%%%%%%%%%%%%%%%%%%%%%%%%%%%%%%%%%%%%%%
This proves (\ref{eq2}) for the case $i\neq j \in \hat{\mathcal{J}}$ and we complete the proof of this lemma.
\end{proof}

\end{document}
